\begin{table}[H]
    \footnotesize
    \centering
    \begin{threeparttable}
        \caption{Novelty-Claim Strength of authors' \(t\)th paper (draft and final)}
        \label{table9llmnovelty}
        \begin{tabular}{p{3cm}S@{}S@{}S@{}S@{}S@{}S@{}}
            \toprule
            &{\(t=1\)}&{\(t=2\)}&{\(t=3\)}&{\(t=4\text{--}5\)}&{\(t\ge6\)}\\
            \midrule
            \multicolumn{6}{l}{\textbf{Predicted \(R_{jP}-R_{jW}\)}}\\
            \quad Women                   &        0.11   &       -0.20   &       -0.44***&       -0.69***&       -0.98***\\
                                          &      (0.21)   &      (0.15)   &      (0.14)   &      (0.20)   &      (0.26)   \\
            \quad Men                     &       -0.31***&       -0.32***&       -0.27***&       -0.22***&       -0.21***\\
                                          &      (0.04)   &      (0.03)   &      (0.02)   &      (0.02)   &      (0.05)   \\
            \midrule\multicolumn{6}{l}{\textbf{Marginal effect of female ratio}}\\
            \quad Published article       &       -0.60** &       -0.58***&       -0.56***&       -0.54***&       -0.51*  \\
                                          &      (0.28)   &      (0.18)   &      (0.13)   &      (0.19)   &      (0.29)   \\
            \quad Draft paper             &       -1.02***&       -0.70***&       -0.39** &       -0.07   &        0.25   \\
                                          &      (0.21)   &      (0.16)   &      (0.19)   &      (0.26)   &      (0.35)   \\
            \midrule
            \textbf{Diff.-in-diff.}&        0.42*  &        0.13   &       -0.17   &       -0.47** &       -0.77***\\
                                          &      (0.24)   &      (0.17)   &      (0.16)   &      (0.21)   &      (0.29)   \\
            \bottomrule
        \end{tabular}
        \begin{tablenotes}
            \tiny
            \item \textit{Notes}. Sample 4,289 observations. Panel one displays magnitude of predicted \(R_{jP}-R_{jW}\) (the direct effect of peer review) for women and men over increasing \(t\). Panel two estimates the marginal effect of an article's female ratio (\(\beta_1+\beta_2\times t\)), separately for draft papers and published articles. Figures from FGLS estimation of~equation~(15), weighted by \(N_{it}\) (see~the data appendix). Control variables include citation count (asinh), \(\text{max. }T\) (author prominence) and \(\text{max. }t\) (author seniority), native speaker and editor and journal-year fixed effects. Standard errors clustered by editor and robust to cross-model correlation in parentheses. ***, ** and * statistically significant at 1\%, 5\% and 10\%, respectively.
        \end{tablenotes}
    \end{threeparttable}
\end{table}